(1) 标量傍轴方程的基模解可通过复光束参量 $q(z)=z+iz_R$ 表示，$z_R=\pi w_0^2/\lambda$。由 $1/q=1/R-i\lambda/(\pi w^2)$ 可得 $w(z)=w_0\sqrt{1+(z/z_R)^2}$；对自由空间传播，$w(z)$ 关于束腰 $z=0$ 对称且取得最小值 $w_0$。
(2) 在 $z=z_R$，$w(z_R)=\sqrt2\,w_0$，束宽增加 $\sqrt2$。远场 $|z|\gg z_R$，$w(z)\approx w_0|z|/z_R$，因而半发散角 $\theta\simeq w_0/z_R=\lambda/(\pi w_0)$（小角近似）。束腰越小，衍射发散越大，这是空间傅里叶带宽的表现。
(3) 波前曲率半径 $R(z)=z[1+(z_R/z)^2]$，束腰处取极限 $R(0)=\infty$，远场时 $R(z)\sim z$。Gouy相位 $\psi(z)=\arctan(z/z_R)$，从负远场到正远场总变化为 $\pi$，在束腰为零。它是傍轴聚焦波的纵向附加相位，不是材料折射率引起的普通相移。
(4) 若 $w_0$ 变为原来的一半，$z_R\propto w_0^2$ 故变为四分之一，远场半发散角 $\theta\propto1/w_0$ 则增大为两倍。采用小束腰可以提高聚焦空间分辨能力，但光束在更短的传播距离内就开始明显发散，准直系统必须权衡两者。
